% Part: applied-modal-logics
% Chapter: epistemic-logic
% Section: truth-at-w

\documentclass[../../../include/open-logic-section]{subfiles}

\begin{document}

\olfileid[pt]{aml}{el}{trw}

\olsection{Verdade em um mundo}

Assim como na lógica modal normal, todo modelo epistêmico determina quais !!{formula}s são verdadeiras em
quais de seus mundos. Usamos a mesma notação ``o modelo $\mModel{M}$ torna a
!!{formula}~$!A$ verdadeira no mundo~$w$'' para a noção básica da semântica
relacional. A relação é definida indutivamente e, para todos os operadores não modais, é idêntica à do caso modal normal.

\begin{defn}\ollabel{defn:mmodels}
  A \emph{verdade de !!a{formula}~$!A$ em~$w$} em um~$\mModel M$, em símbolos
  $\mSat{M}{!A}[w]$, é definida indutivamente da seguinte maneira:
  \begin{enumerate}
  \tagitem{prvFalse}{\indcase{!A}{\lfalse}{Nunca $\mSat{M}{\lfalse}[w]$}.}{}
  \tagitem{prvTrue}{\indcase{!A}{\ltrue}{Sempre $\mSat{M}{\ltrue}[w]$}.}{}
  \item $\mSat{M}{p}[w]$ sse $w \in V(p)$
  \tagitem{prvNot}{\indcase{!A}{\lnot !B}{$\mSat{M}{\indfrm}[w]$ sse
    $\mSat/{M}{!B}[w]$}.}{}
  \tagitem{prvAnd}{\indcase{!A}{(!B \land !C)}{$\mSat{M}{\indfrm}[w]$ sse
    $\mSat{M}{!B}[w]$ e $\mSat{M}{!C}[w]$}.}{}
  \tagitem{prvOr}{\indcase{!A}{(!B \lor !C)}{$\mSat{M}{\indfrm}[w]$ sse
    $\mSat{M}{!B}[w]$ ou $\mSat{M}{!C}[w]$} (ou ambos).}{}
  \tagitem{prvIf}{\indcase{!A}{(!B \lif !C)}{$\mSat{M}{\indfrm}[w]$ sse
    $\mSat/{M}{!B}[w]$ ou $\mSat{M}{!C}[w]$}.}{}
  \tagitem{prvIff}{\indcase{!A}{(!B \liff !C)}{$\mSat{M}{\indfrm}[w]$ sse
    ou ambos $\mSat{M}{!B}[w]$ e $\mSat{M}{!C}[w]$, ou
    nem $\mSat{M}{!B}[w]$ nem $\mSat{M}{!C}[w]$}.}{}
  \item\ollabel{defn:sub:mmodels-box}
    \indcase{!A}{\Knows_a !B}{$\mSat{M}{\indfrm}[w]$ sse
    $\mSat{M}{!B}[w']$ para todo $w' \in W$ tal que $R_a ww'$}
  \end{enumerate} 
\end{defn}

É neste ponto, porém, que precisamos considerar restrições às relações
de acessibilidade. Afinal, pela cláusula~\olref{defn:sub:mmodels-box},
!!a{formula}~$\Knows_a !B$ é verdadeira em~$w$ sempre que não existe~$w'$
tal que $R_a ww'$. Essa é a mesma cláusula da lógica modal normal: quando
um mundo não tem sucessores, todas as fórmulas-$\Box$ são vacuamente
verdadeiras nele. Isso parece extremamente contraintuitivo se
entendermos $\Knows$ como representação de \emph{conhecimento}. Afinal,
tendemos a pensar que não há \emph{nenhuma} circunstância na qual um
agente possa saber ao mesmo tempo tanto $!A$ quanto $\lnot !A$.

Uma solução é garantir que a relação de acessibilidade da lógica
epistêmica seja sempre \emph{reflexiva}. Isso corresponde, em linhas
gerais, à ideia de que o mundo atual é compatível com a informação de
um agente. De fato, as lógicas epistêmicas costumam usar S5, mas outras
podem empregar sistemas mais fracos, dependendo do que exatamente se
pretenda representar pela relação $\Knows_a$.

\begin{figure}
  \begin{center}
    \begin{tikzpicture}[modal]
      \node[world] (w1) [label={[align=right]right:\mTrue{p}\\ \mFalse{q}}]{$w_1$}; 
      \node[world] (w2) [label={[align=right]right:\mTrue{p}\\ \mTrue{q}},
        above right=of w1]{$w_2$}; 
      \node[world] (w3) [label={[align=right]right:\mFalse{p}\\ \mFalse{q}},
        below right=of w1] {$w_3$};
      \draw[<->] (w1) to node {$a$} (w2);
      \draw[->,loop left] (w1) to node {$a,b$} (w1);
      \draw[<->] (w1) to [swap] node {$b$} (w3);
      \draw[->,loop above] (w2) to node {$a,b$} (w2) ;
      \draw[->,loop below] (w3) to node {$a,b$} (w3) ;
    \end{tikzpicture}
  \end{center}
  \caption{Um modelo epistêmico simples.}
  \ollabel{fig:simple}
\end{figure}

\begin{prob}
  Determine quais das afirmações seguintes valem em
  \olref[aml][el][trw]{fig:simple}:
  \begin{enumerate}
    \item $\mSat{M}{\lnot q}[w_1]$;
    \item $\mSat{M}{\Knows_a \lnot q}[w_1]$;
    \item $\mSat{M}{\Knows_b \lnot q}[w_1]$;
    \item $\mSat{M}{\Knows_b q \lor \Knows_b \lnot q}[w_2]$;
    \item $\mSat{M}{\Knows_a( \Knows_b q \lor \Knows_b \lnot q)}[w_2]$;
    \item $\mSat{M}{\EKnows_{\{a,b\}} \lnot q}[w_3]$;
    \end{enumerate}
\end{prob}

Agora que apresentamos nossa definição básica de verdade em um mundo,
os demais conceitos semânticos da lógica modal normal, como validade
modal e consequência, são simplesmente transportados para essa nova
maneira de interpretar os operadores modais.

Também podemos agora formular as condições de verdade do operador de
conhecimento comum~$\CKnows_G$. Recorde, de \olref[sfr][rel][ops]{sec},
que o \emph{fecho transitivo}~$R^+$ de uma relação~$R$ é definido como
\begin{align*}
  R^+ &= \bigcup_{ n \in \mathbb{N}} R^n, 
\intertext{em que}
  R^0 & = R \text{ e}\\ 
  R^{n+1} & = \Setabs{\tuple{x, z}}{\exists y (R^n xy \land Ryz)}.
\end{align*}
Então, se $G$ é um grupo de agentes, definimos $R_G = ( \bigcup_{b
\in G} R_b )^+$ como o fecho transitivo da união das relações de
acessibilidade de todos os agentes.

\begin{defn}
Se $G' \subseteq G$, estabelecemos $\mSat{M}{\CKnows_{G'} !A}[ w]$ sse,
para todo $w'$ tal que $R_{G'} w w'$, $\mSat{M}{!A}[w']$.
\end{defn}

\end{document}
